\subsection{Two explicit obstructions to proposed global shortcuts}

The following refutations concern two proposed auxiliary inferences.
They do not refute D1--D4. In particular, a rational place and a rational
tangential parameter are different requirements.

\begin{claim}{D.20}{Refuted}\label{clm:D20}
The following implication is false: a continuous birational section over
$\mathbf Q$ whose tangential Kummer class is represented by an ordinary
element of $\mathbf Q_p^\times$ at every finite prime, and by a unit at a
set of primes of Dirichlet density one, must have a rational tangential
class or be a unit at all but finitely many primes. The failure occurs
for a section supported at a rational place.
\end{claim}
\begin{proof}
Choose increasing distinct primes $q_j>2^{j^2}$ and define
\[
 A=\prod_{j\geq1}q_j^{j!}\in
 \widehat{\mathbf Q^\times}
 =\varprojlim_n\mathbf Q^\times/(\mathbf Q^\times)^n.
\]
Modulo $n$th powers, every factor with $j\geq n$ is trivial. Thus the
product defines compatible finite-level classes. Its valuation
coordinates are
\[
 v_p(A)=\begin{cases}j!&p=q_j,\\0&p\notin\{q_j:j\geq1\}.
 \end{cases}
\]
Infinitely many are nonzero, so $A$ is not the image of an element of
$\mathbf Q^\times$.

The power completion of $\mathbf Q_p^\times$ is
\[
 \widehat{\mathbf Q_p^\times}
 \simeq p^{\widehat{\mathbf Z}}\times\mathbf Z_p^\times.
\]
Indeed, $\mathbf Q_p^\times=p^{\mathbf Z}\times\mathbf Z_p^\times$,
and the unit group is complete for its finite power quotients. Since
$v_p(A)$ is an ordinary integer, its local image has a unique actual
representative $a_p\in\mathbf Q_p^\times$. One may also construct $a_p$
directly. Remove the one possible factor with $q_j=p$; the remaining
product converges in $\mathbf Z_p^\times$, because $j!$ is eventually
divisible by the exponent of $(\mathbf Z/p^r\mathbf Z)^\times$ for each
$r$. Restore the removed factor. Hence $v_p(a_p)=v_p(A)$.

The set $\{q_j\}$ is infinite but has Dirichlet density zero:
\[
 \sum_j q_j^{-s}\leq\sum_j2^{-j^2}<\infty\qquad(s\geq1),
\]
whereas $\log(1/(s-1))$ tends to infinity as $s\downarrow1$.
The $a_p$ are therefore units at density one, but not cofinitely many,
primes. The real local class is the positive sign class.

We realize these data by a continuous section. Fix a prolongation of
the place $t=0$ of $\mathbf Q(t)$ and its tangential splitting
\[
 D_0\simeq\widehat{\mathbf Z}(1)\rtimes G_{\mathbf Q}
 \subset G_{\mathbf Q(t)}.
\]
For each $j$, choose compatible radicals of $q_j$ and the corresponding
continuous Kummer cocycle
$c_j:G_{\mathbf Q}\to\widehat{\mathbf Z}(1)$. The series
\[
 c(g)=\sum_{j\geq1}j!c_j(g)
\]
converges uniformly in $g$ modulo every $n$, because all terms with
$j\geq n$ vanish modulo $n$. Its finite-level cocycle identities imply
that it is a continuous cocycle. If $s_0$ denotes the standard
tangential section and $i$ the inertia inclusion, then
\[
 s_A(g)=i(c(g))s_0(g)
\]
is a continuous section of $G_{\mathbf Q(t)}\to G_{\mathbf Q}$.
Its pullback of the Kummer class of $t$ is $A$.

For completeness, write a rational function near zero as
$f=t^r u(t)$, where $u$ is regular and nonzero at zero. Hensel lifting
gives compatible roots of $u/u(0)$ in the strict henselian local field,
so the Kummer character of $u$ on $D_0$ is inflated from that of $u(0)$.
Consequently
\[
 s_A^*\kappa(f)=r[A]+\kappa(u(0)).
\]
Thus $s_A^*\kappa(t-c)=\kappa(-c)$ for every nonzero rational $c$,
whereas $s_A^*\kappa(t)=[A]$ is not rational. The underlying rational
place is still exactly zero. The example is cuspidal on
$\mathbf P^1-\{0,1,\infty\}$, and its extra tangential data disappear
in the proper unramified quotient, which kills this inertia. It is
therefore compatible with all the asserted rational-place conclusions.
\end{proof}

\begin{claim}{D.21}{Refuted}\label{clm:D21}
It is false that every fourth-power Fermat twist arising from a section
of $\mathbf P^1_{\mathbf Q}-\{0,1,\infty\}$ admits a nonconstant map
over $\mathbf Q$ to an elliptic curve of rank zero. The twist
\[
 C:\quad x^4+2y^4=3z^4
\]
arises above the rational point $t=1/3$, and every elliptic target of
$C$ over $\mathbf Q$ has positive Mordell--Weil rank.
\end{claim}
\begin{proof}
The plane quartic is smooth, since its three partial derivatives cannot
vanish simultaneously at a projective point; its genus is three. Its
open subset $xyz\ne0$ maps to the thrice-punctured line by
\[
 t=\frac{x^4}{3z^4},\qquad 1-t=\frac{2y^4}{3z^4}.
\]
The rational point $P=(1:1:1)$ lies above $t=1/3$, so this is already
a twist belonging to a geometric section.

The three sign involutions have genus-one quotients. For example, the
$x$-sign quotient has affine equation $v^2=3-2u^4$, where
$u=y/z$ and $v=(x/z)^2$. Its four distinct branch points give genus
one, and the image of $P$ is rational. The other two quotients have the
same description after changing coordinates. Explicit elliptic targets
and maps are
\[
\begin{array}{c|c|c|c}
 &\text{elliptic equation}&X&Y\\\hline
 E_x&Y^2=X^3-6X&-2y^2/z^2&-2x^2y/z^3\\
 E_y&Y^2=X^3-12X&-2x^2/z^2&-4xy^2/z^3\\
 E_z&Y^2=X^3+18X&3x^2/y^2&9xz^2/y^3.
\end{array}
\]
These rational maps extend to the smooth projective curves. To verify
their equations, put $F=x^4+2y^4-3z^4$. After clearing denominators,
$Y^2-X^3-AX$ is respectively $4y^2F$, $8x^2F$, and $-27x^2F$.
Each map factors through its indicated sign quotient. In the quartic
model $v^2=A u^4+B$, the factor map is
$X=A u^2$, $Y=Auv$, with equation $Y^2=X^3+ABX$; its degree is two.
After translating the image of a rational origin, this is an isogeny
of genus-one Jacobians. Thus $E_x,E_y,E_z$ represent the three quotient
isogeny classes.

The images of $P$ are
\[
 (-2,-2)\in E_x(\mathbf Q),\quad
 (-2,-4)\in E_y(\mathbf Q),\quad
 (3,9)\in E_z(\mathbf Q).
\]
They have infinite order. Exhaustive finite-field counts give
\[
\begin{array}{c|cc}
 &\text{first good-reduction count}&\text{second good-reduction count}\\\hline
 E_x&\#E_x(\mathbf F_5)=8&\#E_x(\mathbf F_{13})=18\\
 E_y&\#E_y(\mathbf F_5)=10&\#E_y(\mathbf F_7)=8\\
 E_z&\#E_z(\mathbf F_5)=10&\#E_z(\mathbf F_7)=8.
\end{array}
\]
The discriminant of $Y^2=X^3+AX$ is $-64A^3$, so all displayed
primes are good. Prime-to-$p$ rational torsion injects under reduction
at $p$; see \cite[Chapter~II, Proposition~5.4 and Corollary~5.7]{D_Milne2021}.
For $E_x$ the count at $13$ excludes $5$-primary torsion, and the
count at $5$ excludes $13$-primary torsion. All other primary torsion
orders divide $\gcd(8,18)=2$. The analogous argument at $5,7$ for
$E_y,E_z$ uses $\gcd(10,8)=2$. Each displayed point has nonzero
$Y$-coordinate, so is neither the identity nor of order two. Each
quotient isogeny class has positive rank. The exhaustive calculation is
reproducible from \path{code/D_Fermat_quartic_certificate.py}; its
output is \path{audit/D_Fermat_quartic_certificate.json}.

Finally, let $V$ be the Klein four group of projective sign changes,
with nonidentity elements $\tau_x,\tau_y,\tau_z$. The full quotient
has function field $\mathbf Q(x^2/z^2,y^2/z^2)$ with equation
$X^2+2Y^2=3$. This conic has the rational point $(1,1)$, and hence has
genus zero. On $J=\operatorname{Jac}(C)$ the full norm
$1+\tau_x+\tau_y+\tau_z$ is therefore zero, by pushforward and pullback
for $C\to C/V$. In $\operatorname{End}(J)\otimes\mathbf Q$ the
idempotents
\[
 e_i=\frac{1+\tau_i}{2}\qquad(i=x,y,z)
\]
are pairwise orthogonal and sum to one. Their images are the three
genus-one quotient Jacobians up to isogeny. Hence $J$ is
$\mathbf Q$-isogenous to their product. If $C$ maps nonconstantly to
an elliptic curve $E/\mathbf Q$, translate the image of $P$ to the
origin. The induced map $J\to E$ is nonzero, and so is its restriction
to at least one factor of this product up to isogeny. A nonzero
homomorphism between elliptic curves is an isogeny. Therefore $E$ is
isogenous to one of $E_x,E_y,E_z$ and has positive rank.

The same decomposition excludes a nonconstant map from $C$ to any
abelian variety $A/\mathbf Q$ with $A(\mathbf Q)$ finite: the induced
nonzero homomorphism $J\to A$ is nonzero on some elliptic factor, and
that factor's infinite rational points have infinite image because
its kernel is finite. In particular $C$ cannot map nonconstantly over
$\mathbf Q$ to a positive-genus standard $X_0(N)$ or $X_1(N)$.
Indeed, the proof of Stoll's Corollary~8.8 \cite{D_Stoll2007} supplies
a nonconstant map from each such modular curve to a modular abelian
variety of analytic rank zero, whose rational points are finite.
A nonconstant map of proper curves $C\to M$ is surjective, so its
composition with that map from $M$ would remain nonconstant, a
contradiction. Thus this particular twist does not meet the target
hypotheses of D.22 below.

This excludes a universal rank-zero elliptic-target shortcut on these
twists. It does not exclude additional covers or a different descent
argument, and makes no assertion about the exact ranks or
Tate--Shafarevich groups.
\end{proof}

\paragraph{Why finite localization is not yet finite support.}
Betts--Stix prove that, for a number field containing no CM subfield
and a smooth projective curve of genus at least two, the image of its
Selmer section set in the points over each fixed finite completion is
finite \cite[Theorem~A]{D_BettsStix2025}. The local theorem D.19 supplies
the required point-theoretic localizations at finite places. At a real
place one must separately use the classical real section theorem,
whose target is the set of connected components of the real locus;
the precise input is \cite[Corollary~3.13]{D_Mochizuki2003Real}.
D.19 itself is a $p$-adic theorem and does not address real places.
At a complex place the local Galois group is trivial. Subject to these
archimedean facts, proper global sections are Selmer and the stated
finite-localization theorem applies.

Its conclusion is not the common finite arithmetic support used by
Stix's corollary \cite[Section~3, Corollary~5]{D_Stix2015}. The finite
sets at different completions are not asserted to lie in the base
changes of a single finite subscheme over the number field. Even a
singleton local point can have generic scheme-theoretic image over
the number field if its coordinates are transcendental over that
field. The missing common-support assertion, or the cofinite
integrality condition of D.5, must therefore be supplied by an
additional global argument.

\begin{claim}{D.22}{Proved}\label{clm:D22}
Let $V/k$ be a smooth geometrically connected hyperbolic curve over a
number field, with smooth projective completion $C$. If $C$ admits a
nonconstant morphism to an abelian variety $A/k$ for which $A(k)$ is
finite and $\Sha(k,A)_{\mathrm{div}}=0$, then every continuous section
of $\pi_1(V)\to G_k$ is a rational-point section or is cuspidal.

In particular this conclusion holds over $\mathbf Q$ whenever $C$
admits a nonconstant $\mathbf Q$-morphism to a standard modular curve
$X_0(N)$ or $X_1(N)$ of positive genus. It includes the proper modular
curves of genus at least two and hyperbolic open modular curves with
positive-genus completions.
\end{claim}
\begin{proof}
Let $s$ be a continuous section. At every finite place $v$, D.6, using
the proper local theorem D.19, makes $s_v$ a point or cusp section.
Its image in the fundamental group of $C$ is therefore induced by a
point $x_v\in C(k_v)$. At every real place, the profinite real section
theorem \cite[Corollary~3.13]{D_Mochizuki2003Real} identifies sections
for the hyperbolic curve $V$ with connected components of $V(\mathbf R)$.
Choose a point in that component and map it to $C$. This specifies the
required real component of the modified adelic point; it does not
assert that a real point within a component is uniquely determined.
Complex factors are omitted in the corrected modified adelic space
\cite[Section~1]{D_StollErrata2017}. Thus $s$ supplies
\[
 (x_v)\in C(\mathbf A_k)_\bullet^{\mathrm{f\text{-}cov}}
 \subseteq C(\mathbf A_k)_\bullet^{\mathrm{f\text{-}ab}}.
\]
Here the full finite-descent assertion follows from Harari--Stix,
Theorem~2.1 and Remark~2.2(1) in the supplied preprint
\cite{D_HarariStix2012}: each finite torsor's evaluations are obtained
by restriction of its pullback along the global section. This is a
condition on torsors under arbitrary finite group schemes, not only
constant or geometrically connected torsors.

Under the stated abelian-target hypothesis, Stoll's Theorem~8.6
\cite{D_Stoll2007} gives
\[
 C(\mathbf A_k)_\bullet^{\mathrm{f\text{-}ab}}=C(k).
\]
Consequently the finite-place local points of $s$ all equal one
$x\in C(k)$, and their real components equal those of $x$. Choosing
$x$ as the real representative gives common finite support $\{x\}$.
If $s$ is cuspidal there is nothing further to prove. Otherwise Stix's
finite-support criterion, Corollary~5 in \cite{D_Stix2015}, shows
that $s$ is the section of a point of $V(k)$; its image in $C$ is
necessarily $x$. For a proper hyperbolic curve this is the point-section
conclusion, since no cusps occur.

For the modular assertion put $M=X_0(N)$ or $X_1(N)$, with
$g(M)>0$, and let $f:C\to M$ be nonconstant. Stoll's Corollary~8.8
states precisely that
$M(\mathbf A_{\mathbf Q})_\bullet^{\mathrm{f\text{-}ab}}=M(\mathbf Q)$.
Functoriality sends $(x_v)$ to one $y\in M(\mathbf Q)$.
The scheme $Z=f^{-1}(y)$ is finite over $\mathbf Q$, and
\[
 x_v\in Z(\mathbf Q_v)\qquad\text{for every finite place }v.
\]
Since $g(C)\geq g(M)>0$, Stoll's Theorem~8.2 applies to $Z\subset C$
and makes $(x_v)$ a diagonal rational point of $C$. This is the
finite-fibre proof of Proposition~8.5, with exactly the repair in
\cite[Section~3]{D_StollErrata2017}: no assertion that the specified
real component meets $Z(\mathbf R)$ is needed before applying
Theorem~8.2; a density-one set of finite places suffices. The
finite-support argument of the preceding paragraph now applies.
\end{proof}

\paragraph{Consequence for selected lifts of sections of $U$.}
If a section of $U=\mathbf P^1_{\mathbf Q}-\{0,1,\infty\}$ lifts to a
finite \'etale cover $V\to U$ whose completion satisfies D.22, its lift
is a point or cusp section. A finite \'etale map pushes such sections
to point or cusp sections, respectively: in the finite map between
the smooth completions the complement of $V$ lies over the complement
of $U$. This settles the original section in that case. Arbitrary
arithmetic twists of covers of $U$ are not thereby shown to map over
$\mathbf Q$ to a standard positive-genus modular curve or to a
rank-zero abelian target. Establishing that property for the remaining
twist classes would be an additional theorem. D.22 does not assert
universal D3 or D4.
